Dane testowe zostały wygenerowane przy pomocy autorskiego skryptu napisanego
w języku Python, zlokalizowanego w module \texttt{src/generators/data\_generator.py}.
Generator korzysta z biblioteki Faker w wersji obsługującej lokalizacje
\texttt{pl\_PL} oraz \texttt{en\_US}, co umożliwia tworzenie realistycznych
danych osobowych z mieszanką polskich i anglojęzycznych imion, nazwisk oraz
opisów tekstowych. Wszystkie generowane wartości pseudolosowe są
deterministycznie powtarzalne dzięki ustawieniu wspólnego ziarna (\texttt{seed=42})
zarówno dla biblioteki Faker, jak i dla wbudowanego generatora liczb
losowych Pythona, co gwarantuje powtarzalność wyników testów wydajnościowych
pomiędzy uruchomieniami.

Konfiguracja wolumenu danych została zdefiniowana w module \texttt{src/config.py}
w postaci trzech profili: \textit{small}, \textit{medium} oraz \textit{large}.
Profil small obejmuje 5\,000 użytkowników, 2\,000 osób w katalogu twórców,
2\,000 tytułów oraz 300\,000 sesji oglądania. Profil medium zawiera
odpowiednio 20\,000 użytkowników, 5\,000 osób, 5\,000 tytułów oraz 500\,000
sesji oglądania. Profil large powiększa zbiór do 100\,000 użytkowników,
20\,000 osób, 15\,000 tytułów oraz 6\,000\,000 sesji oglądania. Łączna liczba
rekordów we wszystkich tabelach docelowo osiąga około pół miliona, miliona
oraz dziesięciu milionów wierszy odpowiednio dla wolumenów small, medium
oraz large, co umożliwia obserwację skalowania charakterystyk wydajnościowych
silników względem rozmiaru zbioru.

W pierwszym etapie generacji danych tworzone są encje
niezależne -- użytkownicy, osoby z katalogu twórców oraz treści. Następnie
generowane są encje powiązane jeden-do-wielu z użytkownikiem (profile,
subskrypcje oraz płatności), encje powiązane z treścią (powiązania
z osobami, sezony, odcinki) oraz, w ostatnim etapie, rekordy aktywności
profili (historia oglądania, lista do obejrzenia, oceny). Każda tabela
zapisywana jest w postaci osobnego pliku CSV z kodowaniem UTF-8, co
umożliwia jednolity import do wszystkich czterech systemów bazodanowych.

Procedura generowania użytkowników wprowadza realistyczne rozkłady
prawdopodobieństwa odzwierciedlające typową dystrybucję rynkową. Kraj
zamieszkania jest losowany z wagą 40\% dla Polski, 15\% dla Stanów
Zjednoczonych oraz mniejszymi udziałami pozostałych ośmiu krajów
europejskich. Numer telefonu jest poprzedzany prefiksem zgodnym z krajem
i jest obecny u około 70\% użytkowników. Status konta jest losowany
z rozkładu, w którym użytkownicy aktywni stanowią dominującą większość,
przy czym uwzględniono również kategorie \textit{inactive}, \textit{suspended}
oraz \textit{deleted}, co pozwala na realistyczne modelowanie zapytań
filtrujących po stanie konta. Adresy e-mail konstruowane są na podstawie
imienia, nazwiska oraz identyfikatora użytkownika, z wykorzystaniem zbioru
typowych domen polskich i międzynarodowych. Wartość pola
\texttt{password\_hash} jest pseudolosowym ciągiem znaków zgodnym
z formatem skrótu bcrypt.

Liczba profili przypadających na konto jest losowana z rozkładu, w którym
najczęstsza wartość mieści się w przedziale od dwóch do trzech profili,
co odpowiada typowemu wzorcowi użycia rodzinnego. Profil drugi w obrębie
konta jest oznaczany jako profil dziecięcy z prawdopodobieństwem 40\%,
co automatycznie ustawia jego klasyfikację wiekową na wartość \textit{ALL}.
Liczba subskrypcji per użytkownik również jest losowana z odpowiedniego
rozkładu, przy czym ostatnia subskrypcja zawsze posiada status
\textit{active}, natomiast wcześniejsze są oznaczane jako \textit{cancelled}
lub \textit{expired}, co umożliwia testowanie zapytań analizujących historię
zmian planu. Plan subskrypcji wybierany jest spośród trzech wariantów --
\textit{Basic}, \textit{Standard} oraz \textit{Premium} -- z wagami 30\%,
45\% oraz 25\% odpowiednio.

Treści są dzielone na filmy oraz seriale w stosunku 60\% do 40\%. Filmy
posiadają określoną długość, natomiast seriale otrzymują strukturę
podrzędną w postaci sezonów oraz odcinków. Liczba sezonów na serial
podlega rozkładowi malejącemu, w którym dominują seriale jedno-,
dwu- oraz trzysezonowe. Liczba odcinków w sezonie jest losowana
z przedziału od czterech do szesnastu, z datami premier rozłożonymi
w cotygodniowych odstępach. Każdy tytuł otrzymuje od jednej do trzech
etykiet gatunkowych losowanych z dwunastoelementowego zbioru.

Dla każdego tytułu generowany jest również obiekt JSON wypełniający
pole \texttt{metadata}. Obiekt zawiera nazwę studia produkcyjnego,
budżet w przedziale od 500\,000 do 200\,000\,000 jednostek, listę
otrzymanych nagród (od zera do trzech pozycji), listę znaczników
kategoryzujących, listę krajów koprodukcji oraz zagnieżdżony obiekt
\texttt{streaming\_quality} z parametrami technicznymi strumienia.
Tak skonstruowane pole umożliwia testowanie scenariuszy odpytywania
danych zagnieżdżonych w sposób reprezentatywny dla rzeczywistych
zastosowań platform streamingowych.

Powiązania między treściami a osobami są generowane z zachowaniem
struktury produkcyjnej -- każdy tytuł otrzymuje dokładnie jednego
reżysera, od jednego do dwóch scenarzystów oraz od trzech do dwunastu
aktorów, każdy z przypisaną nazwą postaci oraz pozycją w napisach.
Sesje oglądania są przypisywane losowym profilom oraz losowym treściom,
przy czym dla seriali wybierany jest losowy odcinek z dostępnych. Postęp
oglądania jest losowany z rozkładu jednostajnego od zera do stu procent,
a wartości przekraczające próg 95\% są normalizowane do stu procent
i oznaczane znacznikiem ukończenia. Oceny wystawiane przez profile
podlegają rozkładowi pochylonemu w stronę wartości pozytywnych
(z modą w przedziale od siedmiu do ośmiu punktów), co odzwierciedla
typową asymetrię ocen obserwowaną na platformach streamingowych.
Tekst recenzji jest dołączany w 30\% przypadków.

Wygenerowane pliki CSV stanowią wspólny punkt wyjścia dla mechanizmów
ładowania danych do każdego z czterech systemów bazodanowych. Konwersja
do modelu dokumentowego MongoDB realizowana jest przez dedykowany skrypt
importujący, który łączy plik \texttt{users.csv} z plikami \texttt{profiles.csv}
oraz \texttt{subscriptions.csv} w pojedynczy dokument oraz osadza dane
gatunków, obsady, sezonów i odcinków wewnątrz dokumentów kolekcji
\texttt{content}. Konwersja do modelu grafowego Neo4j realizowana jest
narzędziem \texttt{neo4j-admin import}, które przyjmuje pliki CSV
opisujące węzły oraz relacje w formacie zgodnym ze specyfikacją silnika.
Mechanizmy ładowania danych dla poszczególnych silników zostały szczegółowo
opisane w dalszej części pracy.
